% ============================================================================
% CAPÍTULO 7: DISCUSIÓN
% ============================================================================

\chapter{Discusión}
\label{cap:discusion}

El análisis de los resultados obtenidos en condiciones de vida libre (free-living) bajo el enfoque \textit{Bring Your Own Device} (BYOD) aporta evidencia valiosa sobre la factibilidad y validez del uso de dispositivos portátiles comerciales para la evaluación del comportamiento sedentario. Este paradigma metodológico representa una transición significativa en la investigación biomédica, al desplazar la observación controlada hacia contextos ecológicamente válidos donde los individuos realizan sus actividades cotidianas sin interferencia experimental. Tal desplazamiento no solo amplía la representatividad de los datos, sino que también facilita la generación de conocimiento aplicable a entornos reales de salud pública \cite{Doherty2021,Liu2022}.

En los últimos años, los estudios basados en wearables han demostrado una creciente capacidad para cuantificar la actividad física y la variabilidad de la frecuencia cardiaca con precisión comparable a los métodos de laboratorio \cite{Strain2020,Henriksen2018}. Sin embargo, la literatura actual continúa señalando brechas en la interpretación de estos datos, particularmente cuando se pretende traducir métricas continuas en información clínicamente relevante. En este sentido, el presente modelo difuso contribuye a superar una de las limitaciones centrales de los estudios BYOD: la heterogeneidad de las señales fisiológicas y la dificultad de integrar dimensiones objetivas (biométricas) y subjetivas (percepción de bienestar) en un marco analítico coherente.

El sistema de inferencia difusa propuesto permite representar relaciones no lineales entre las variables de sedentarismo, variabilidad autonómica (HRV-SDNN) y percepción de calidad de vida, conservando la interpretabilidad lingüística necesaria para la toma de decisiones clínicas y de salud pública. A diferencia de los algoritmos de caja negra, el modelo Mamdani facilita la comprensión del proceso decisional mediante reglas semánticas (``si actividad es baja y HRV es baja, entonces riesgo es alto''), lo que habilita su implementación en plataformas digitales de monitoreo comunitario y programas de intervención personalizados \cite{Kaur2022,Escalante2023}.

Desde la perspectiva de la salud poblacional, los resultados de este estudio refuerzan la noción de que la evaluación continua del sedentarismo no debe limitarse al recuento de pasos o al tiempo total sentado, sino que debe considerar la dinámica temporal y fisiológica del comportamiento. Los patrones observados en condiciones de vida libre evidencian que las interrupciones breves del sedentarismo, aun sin alcanzar los criterios clásicos de ejercicio, pueden asociarse con una mejora del tono autonómico y de la percepción subjetiva de vitalidad. Este hallazgo coincide con los planteamientos del \textit{Global Action Plan on Physical Activity 2018--2030} de la OMS, que enfatiza la importancia de reducir los periodos prolongados de inactividad más allá del volumen total de actividad física \cite{WHO2018,Bull2020}.

Asimismo, la metodología BYOD ofrece ventajas logísticas y económicas para su escalamiento a gran escala, dado que aprovecha la infraestructura tecnológica personal existente. Esto permite el monitoreo remoto de amplias poblaciones sin requerir equipamiento especializado o visitas presenciales, reduciendo costos y favoreciendo la equidad en el acceso a la investigación y al seguimiento clínico. En este contexto, la presente investigación sienta las bases para el desarrollo de modelos predictivos accesibles, basados en datos pasivos de wearables, que podrían integrarse a programas de prevención de enfermedades crónicas o estrategias nacionales de vigilancia epidemiológica del comportamiento sedentario.

La aplicabilidad práctica de estos resultados radica en su potencial para orientar políticas públicas centradas en el uso responsable de la tecnología digital como herramienta de salud. Los modelos explicables, como el propuesto, permiten diseñar dashboards y alertas personalizadas que incentiven la movilidad cotidiana, identifiquen patrones de riesgo y faciliten la toma de decisiones en entornos laborales, educativos y comunitarios. Su naturaleza adaptativa posibilita, además, la personalización de intervenciones conductuales, incorporando recomendaciones dinámicas según la respuesta fisiológica de cada individuo.

Finalmente, este estudio demuestra que la convergencia entre IA explicable, dispositivos portátiles y estrategias de salud pública constituye un eje transformador en la vigilancia de los comportamientos sedentarios. La combinación de metodologías BYOD con modelos difusos no solo amplía la comprensión del fenómeno desde la perspectiva fisiológica y conductual, sino que también abre el camino hacia ecosistemas de salud digital basados en datos de vida real, donde la prevención y la promoción de la salud se integren de forma continua y personalizada.

En conjunto, los hallazgos aquí presentados fortalecen la premisa de que la medición objetiva del sedentarismo mediante tecnologías portátiles, interpretada a través de sistemas inteligentes transparentes, puede convertirse en una herramienta estratégica para la toma de decisiones en salud pública, contribuyendo a reducir la carga de enfermedad asociada a la inactividad y a promover estilos de vida más activos y sostenibles.

% ============================================================================
% SECCIONES OPCIONALES - COMENTADAS (PLANTILLAS ILUSTRATIVAS)
% ============================================================================
%
% Las siguientes secciones son plantillas ilustrativas que pueden activarse
% según las necesidades específicas del documento
%
% % ----------------------------------------------------------------------------
% % SECCIÓN 7.1: Discusión de Hallazgos Principales
% % ----------------------------------------------------------------------------
% \section{Discusión de Hallazgos Principales}
% \label{sec:discusion_hallazgos}
%
% \subsection{[Hallazgo Principal 1]}
% \label{subsec:discusion_hallazgo1}

% ESTRUCTURA SUGERIDA PARA CADA HALLAZGO:
% 1. Recuerda brevemente el hallazgo
% 2. Interpreta: ¿qué significa?
% 3. Compara con la literatura
% 4. Explica similitudes/diferencias
% 5. Proporciona explicaciones teóricas
%
% \textbf{Hallazgo:}
%
% En este estudio se encontró que \textit{[resume brevemente el hallazgo principal 1]}.
%
% \textbf{Interpretación:}
%
% Este hallazgo sugiere que \textit{[tu interpretación del significado]}, lo cual implica \textit{[implicaciones]}.
%
% \textbf{Comparación con la literatura:}
%
% Este resultado es \textit{[consistente/inconsistente]} con lo reportado por \cite{autor1}, quien encontró que \textit{[qué encontró]}. Similarmente, \cite{autor2} reportó \textit{[hallazgo similar/diferente]}.
%
% Sin embargo, difiere de los hallazgos de \cite{autor3}, quienes reportaron \textit{[diferencia]}. Esta discrepancia podría explicarse por \textit{[posibles razones: diferencias en población, metodología, contexto, etc.]}.
%
% \textbf{Explicación teórica:}
%
% Desde la perspectiva de \textit{[teoría o modelo del marco teórico]}, este resultado se explica porque \textit{[explicación basada en teoría]} \cite{referencia_teoria}.
%
% \subsection{[Hallazgo Principal 2]}
% \label{subsec:discusion_hallazgo2}
%
% [Repite la estructura para cada hallazgo principal]
%
% \subsection{[Hallazgo Principal 3]}
% \label{subsec:discusion_hallazgo3}
%
% [Continúa con todos tus hallazgos principales]
%
% % ----------------------------------------------------------------------------
% % SECCIÓN 7.2: Integración de Resultados con el Cuadro Comparativo
% % ----------------------------------------------------------------------------
% \section{Comparación con Investigaciones Previas}
% \label{sec:comparacion_investigaciones}
%
% % INSTRUCCIONES:
% % - Retoma el cuadro comparativo del Cap. 2 (Marco Teórico)
% % - Discute cómo tus resultados se comparan con esos estudios
%
% \subsection{Síntesis Comparativa}
% \label{subsec:sintesis_comparativa}
%
% La \Cref{tab:sintesis_discusion} presenta una síntesis comparativa entre los hallazgos de estudios previos y los resultados de esta investigación.
%
% \begin{table}[htbp]
%     \centering
%     \caption{Comparación de hallazgos principales con investigaciones previas}
%     \label{tab:sintesis_discusion}
%     \small
%     \begin{tabular}{@{}p{3cm}p{4cm}p{4cm}p{3cm}@{}}
%         \toprule
%         \textbf{Estudio} & \textbf{Hallazgo Principal} & \textbf{Hallazgo de Esta Investigación} & \textbf{Consistencia} \\
%         \midrule
%         
%         Autor1 et al. (2023) & 
%         [Hallazgo del estudio] & 
%         [Tu hallazgo relacionado] & 
%         \textit{Consistente/ Parcialmente/ Inconsistente} \\
%         \midrule
%         
%         Autor2 \& Autor3 (2022) & 
%         [Hallazgo] & 
%         [Tu hallazgo] & 
%         \textit{[Consistencia]} \\
%         \midrule
%         
%         [Continúa con más estudios] & 
%         ... & 
%         ... & 
%         ... \\
%         
%         \bottomrule
%     \end{tabular}
% \end{table}
%
% \subsection{Análisis de Convergencias}
% \label{subsec:convergencias}
%
% Los resultados de esta investigación convergen con la literatura previa en los siguientes aspectos:
%
% \begin{enumerate}
%     \item \textbf{Aspecto 1:} Al igual que \cite{autor1} y \cite{autor2}, este estudio confirma que \textit{[punto de convergencia]}. Esto sugiere que \textit{[implicación]}.
%     
%     \item \textbf{Aspecto 2:} \textit{[Otra convergencia]}...
%     
%     \item \textbf{Aspecto 3:} \textit{[Otra convergencia]}...
% \end{enumerate}
%
% \subsection{Análisis de Divergencias}
% \label{subsec:divergencias}
%
% Por otro lado, se identificaron diferencias con algunos estudios previos:
%
% \begin{enumerate}
%     \item \textbf{Diferencia 1:} A diferencia de \cite{autor3}, quien reportó \textit{[hallazgo diferente]}, en este estudio se encontró \textit{[tu hallazgo]}. 
%     
%     \textbf{Posibles explicaciones:}
%     \begin{itemize}
%         \item Diferencias en la población estudiada: \textit{[explica]}
%         \item Diferencias metodológicas: \textit{[explica]}
%         \item Diferencias contextuales: \textit{[explica]}
%     \end{itemize}
%     
%     \item \textbf{Diferencia 2:} \textit{[Otra divergencia con explicación]}...
% \end{enumerate}
%
% % ----------------------------------------------------------------------------
% % SECCIÓN 7.3: Logros de la Investigación
% % ----------------------------------------------------------------------------
% \section{Logros de la Investigación}
% \label{sec:logros}
%
% Esta investigación logró:
%
% \begin{enumerate}
%     \item \textbf{Logro 1:} Se cumplió el objetivo de \textit{[objetivo]} al \textit{[cómo se logró]}, lo cual representa una aportación importante porque \textit{[por qué es importante]}.
%     
%     \item \textbf{Logro 2:} Se confirmó la hipótesis de que \textit{[hipótesis]}, con evidencia estadística significativa ($p < $ \textit{[valor]}), lo cual \textit{[implicación]}.
%     
%     \item \textbf{Logro 3:} Se logró \textit{[validar/desarrollar/adaptar]} el instrumento \textit{[nombre]} en el contexto de \textit{[población]}, con adecuadas propiedades psicométricas ($\alpha$ = \textit{[valor]}), lo cual contribuye metodológicamente al campo.
%     
%     \item \textbf{Logro 4:} Se identificaron predictores significativos de \textit{[variable]}, específicamente \textit{[predictores]}, explicando el \textit{[\%]} de la varianza, lo cual tiene implicaciones prácticas para \textit{[aplicación]}.
%     
%     \item \textbf{Logro 5:} Se generó conocimiento nuevo sobre \textit{[aspecto específico]} que no había sido explorado previamente en \textit{[contexto/población]}.
% \end{enumerate}
%
% % ----------------------------------------------------------------------------
% % SECCIÓN 7.4: Limitaciones del Estudio
% % ----------------------------------------------------------------------------
% \section{Limitaciones del Estudio}
% \label{sec:limitaciones}
%
% % INSTRUCCIONES:
% % - Sé HONESTO sobre las limitaciones
% % - Esto demuestra rigor científico, NO debilita tu trabajo
% % - Categoriza las limitaciones (metodológicas, contextuales, muestrales, etc.)
%
% A pesar de los logros alcanzados, es importante reconocer las siguientes limitaciones:
%
% \subsection{Limitaciones Metodológicas}
% \label{subsec:limitaciones_metodologicas}
%
% \begin{enumerate}
%     \item \textbf{Diseño del estudio:} El diseño \textit{[transversal/observacional/etc.]} no permite establecer relaciones causales, únicamente asociaciones. Para determinar causalidad se requeriría un diseño \textit{[experimental/longitudinal]}.
%     
%     \item \textbf{Instrumentos:} La recolección de datos se basó en \textit{[autorreporte/cuestionarios]}, lo cual puede estar sujeto a sesgos de \textit{[deseabilidad social/memoria/etc.]}. Mediciones objetivas (como \textit{[ejemplos]}) podrían complementar los hallazgos.
%     
%     \item \textbf{Variables confusoras:} Aunque se controlaron \textit{[variables controladas]}, es posible que otras variables no medidas (como \textit{[ejemplos]}) puedan estar influyendo en los resultados.
% \end{enumerate}
%
% \subsection{Limitaciones Muestrales}
% \label{subsec:limitaciones_muestrales}
%
% \begin{enumerate}
%     \item \textbf{Tamaño de muestra:} Si bien la muestra fue suficiente para detectar efectos significativos, un tamaño muestral mayor ($n > $ \textit{[número]}) podría permitir análisis de subgrupos y aumentar el poder estadístico.
%     
%     \item \textbf{Tipo de muestreo:} El muestreo \textit{[tipo]} puede limitar la generalización de los resultados a \textit{[población más amplia]}. Estudios futuros con muestreo probabilístico fortalecerían la validez externa.
%     
%     \item \textbf{Características de la muestra:} La muestra estuvo compuesta principalmente por \textit{[característica, ej: adultos jóvenes, mujeres, etc.]}, lo cual puede limitar la aplicabilidad a \textit{[otros grupos]}.
% \end{enumerate}
%
% \subsection{Limitaciones Contextuales}
% \label{subsec:limitaciones_contextuales}
%
% \begin{enumerate}
%     \item \textbf{Contexto geográfico:} Los resultados provienen de \textit{[Chihuahua/contexto específico]}, por lo que la generalización a otras regiones con características \textit{[sociodemográficas/culturales/etc.]} diferentes debe hacerse con cautela.
%     
%     \item \textbf{Período temporal:} Los datos se recolectaron durante \textit{[período]}, lo cual podría haber sido influenciado por \textit{[eventos contextuales como pandemia, estacionalidad, etc.]}.
% \end{enumerate}
%
% \subsection{Limitaciones de Recursos}
% \label{subsec:limitaciones_recursos}
%
% \begin{enumerate}
%     \item \textbf{Recursos económicos:} Los recursos disponibles limitaron \textit{[aspecto, ej: tamaño de muestra, número de mediciones, uso de tecnología específica]}.
%     
%     \item \textbf{Tiempo:} El tiempo disponible para el estudio (\textit{[duración]}) no permitió \textit{[seguimiento longitudinal/mediciones repetidas/etc.]}.
% \end{enumerate}
%
% % ----------------------------------------------------------------------------
% % SECCIÓN 7.5: Riesgos Identificados
% % ----------------------------------------------------------------------------
% \section{Riesgos Identificados}
% \label{sec:riesgos}
%
% Durante el desarrollo de esta investigación se identificaron los siguientes riesgos:
%
% \begin{enumerate}
%     \item \textbf{Sesgo de selección:} \textit{[Describe el riesgo y cómo se mitigó]}
%     
%     \item \textbf{Sesgo de respuesta:} \textit{[Describe el riesgo y cómo se mitigó]}
%     
%     \item \textbf{Pérdida de participantes:} \textit{[Si aplica, describe la tasa de pérdida y su posible impacto]}
% \end{enumerate}
%
% % ----------------------------------------------------------------------------
% % SECCIÓN 7.6: Implicaciones de los Hallazgos
% % ----------------------------------------------------------------------------
% \section{Implicaciones de los Hallazgos}
% \label{sec:implicaciones}
%
% \subsection{Implicaciones Teóricas}
% \label{subsec:implicaciones_teoricas}
%
% Los resultados de esta investigación tienen las siguientes implicaciones teóricas:
%
% \begin{enumerate}
%     \item Apoyan la teoría de \textit{[teoría]} al confirmar que \textit{[aspecto]}
%     \item Sugieren una revisión de \textit{[concepto/modelo]} en el sentido de \textit{[sugerencia]}
%     \item Aportan evidencia sobre \textit{[mecanismo/proceso]} que no había sido claramente demostrado
% \end{enumerate}
%
% \subsection{Implicaciones Prácticas}
% \label{subsec:implicaciones_practicas}
%
% Los hallazgos tienen aplicaciones prácticas inmediatas:
%
% \begin{enumerate}
%     \item \textbf{Para profesionales de la salud:} Los resultados sugieren que \textit{[recomendación práctica específica]}
%     
%     \item \textbf{Para tomadores de decisiones:} Se recomienda \textit{[política/programa]} basándose en \textit{[hallazgo]}
%     
%     \item \textbf{Para la población objetivo:} \textit{[Implicación para beneficiarios]}
% \end{enumerate}
%
% \subsection{Implicaciones Metodológicas}
% \label{subsec:implicaciones_metodologicas}
%
% Metodológicamente, este estudio aporta:
%
% \begin{enumerate}
%     \item Un instrumento validado para \textit{[propósito]} en \textit{[población]}
%     \item Una estrategia metodológica para \textit{[proceso]}
%     \item Evidencia sobre la utilidad de \textit{[técnica/enfoque]}
% \end{enumerate}
%
% % ----------------------------------------------------------------------------
% % SECCIÓN 7.7: Reflexión sobre el Cumplimiento de Objetivos e Hipótesis
% % ----------------------------------------------------------------------------
% \section{Cumplimiento de Objetivos e Hipótesis}
% \label{sec:cumplimiento}
%
% \subsection{Objetivo General}
% \label{subsec:cumplimiento_og}
%
% El objetivo general de \textit{[enunciado]} fue alcanzado satisfactoriamente al \textit{[cómo se logró]}.
%
% \subsection{Objetivos Específicos}
% \label{subsec:cumplimiento_oe}
%
% \begin{itemize}
%     \item \textbf{OE1:} \textit{[Enunciado]} - \textbf{Cumplido.} \textit{[Evidencia del cumplimiento]}
%     \item \textbf{OE2:} \textit{[Enunciado]} - \textbf{Cumplido.} \textit{[Evidencia]}
%     \item \textbf{OE3:} \textit{[Enunciado]} - \textbf{Cumplido.} \textit{[Evidencia]}
%     \item \textbf{OE4:} \textit{[Enunciado]} - \textbf{Cumplido.} \textit{[Evidencia]}
% \end{itemize}
%
% \subsection{Hipótesis}
% \label{subsec:cumplimiento_hipotesis}
%
% \textbf{Hipótesis de Investigación (Hi):} \textit{[Enunciado]}
%
% \textbf{Resultado:} La hipótesis fue \textit{[confirmada/rechazada]} con base en \textit{[evidencia estadística: prueba utilizada, valor estadístico, p-value]}.
%
% \textit{[Interpreta qué significa este resultado]}

% ============================================================================
% V2 - DISCUSIÓN TEÓRICA DETALLADA (PROPUESTA ALTERNATIVA)
% ============================================================================
\section{V2\_Discusión}
\label{sec:v2_discusion}

El hallazgo principal de este estudio indica que el modelo de evaluación del comportamiento sedentario basado en lógica difusa logró integrar de manera robusta variables biométricas obtenidas en condiciones de vida libre mediante dispositivos \textit{wearables} de uso personal bajo el enfoque \textit{Bring Your Own Device} (BYOD). La relación identificada entre la variabilidad de la frecuencia cardiaca (HRV-SDNN), los patrones horarios de movimiento y la percepción de calidad de vida relacionada con la salud (CVRS) revela que la estabilidad autonómica se asocia con menores niveles de sedentarismo acumulado y con una percepción más favorable de vitalidad y funcionamiento general. Este hallazgo sugiere que las métricas fisiológicas obtenidas de manera pasiva, cuando son interpretadas a través de un sistema difuso explicable, pueden constituir indicadores fiables del equilibrio entre carga y recuperación en la vida cotidiana. Lo anterior implica una oportunidad concreta para redefinir la vigilancia de la salud desde una perspectiva continua, dinámica y personalizada.

Estos resultados son consistentes con lo reportado por Doherty et al. (2021) y Strain et al. (2020), quienes demostraron que los sensores portátiles pueden registrar con alta precisión el gasto energético y los niveles de actividad física en contextos no controlados, así como predecir el riesgo futuro de enfermedad cardiovascular. De manera similar, Henriksen et al. (2018) confirmaron la utilidad de los dispositivos comerciales para la evaluación de patrones de actividad en estudios longitudinales. Sin embargo, a diferencia de los modelos basados en regresión o aprendizaje profundo empleados en tales investigaciones, el presente modelo difuso prioriza la interpretabilidad, permitiendo traducir el comportamiento fisiológico en reglas lingüísticas comprensibles para el usuario y el profesional de la salud. Esta diferencia metodológica podría explicar la mayor aplicabilidad práctica del modelo propuesto en contextos comunitarios y clínicos, donde la transparencia algorítmica es un requisito ético y operativo.

Desde la perspectiva teórica, el resultado puede interpretarse dentro del marco de la teoría de la autorregulación fisiológica y el modelo de carga alostática, según los cuales el organismo mantiene un equilibrio dinámico entre estrés y recuperación. Una disminución sostenida en la HRV reflejaría un estado de sobrecarga autonómica que, al combinarse con largos periodos de inactividad, incrementa el riesgo de deterioro metabólico y psicológico. El modelo difuso reproduce este razonamiento en términos cuantitativos, asignando grados intermedios de pertenencia a los estados ``activo'', ``moderadamente activo'' o ``sedentario'', con lo cual captura la naturaleza gradual del comportamiento humano en la vida real. Este enfoque explica por qué los resultados se mantienen coherentes incluso ante la variabilidad interindividual y la ausencia de condiciones de laboratorio, confirmando la capacidad del sistema para adaptarse a entornos BYOD y a la heterogeneidad tecnológica de los dispositivos.

En comparación con investigaciones previas, los resultados de este trabajo convergen en tres aspectos esenciales. Primero, coinciden con los reportes de Bull et al. (2020) y la OMS (2018) al destacar que la reducción de los periodos prolongados de sedentarismo tiene un impacto independiente sobre la salud, aun cuando no se incremente el volumen total de actividad física. Segundo, refuerzan la observación de Liu y Chen (2022), quienes señalaron que los enfoques BYOD amplían la representatividad de la población al incluir datos capturados en escenarios cotidianos. Tercero, se alinean con las recomendaciones de Escalante et al. (2023) respecto al uso de inteligencia artificial explicable para garantizar la transparencia y la trazabilidad de los algoritmos aplicados a datos personales de salud. Estas convergencias sustentan la validez externa y el potencial ético-metodológico del presente modelo.

No obstante, se identifican divergencias relevantes con ciertos trabajos que se apoyan exclusivamente en métricas agregadas de pasos o tiempo total sentado, como los de Shamah-Levy et al. (2023), donde no se consideró la dimensión autonómica de la regulación fisiológica. En este estudio, la inclusión del HRV-SDNN permitió detectar patrones de sedentarismo con sensibilidad al estrés fisiológico, lo cual genera una lectura más integral del bienestar. Esta discrepancia puede atribuirse a diferencias metodológicas en la resolución temporal de los datos y a la omisión de la variabilidad cardiaca como marcador de fatiga o carga interna. En consecuencia, el presente trabajo amplía el marco interpretativo del sedentarismo al incorporar indicadores fisiológicos de adaptación.

Los logros de esta investigación se centran en tres niveles. En primer lugar, se alcanzó el objetivo general de diseñar y validar un modelo difuso explicable capaz de evaluar el comportamiento sedentario a partir de datos de vida libre, cumpliendo así el principio de aplicabilidad ecológica. En segundo lugar, se demostró que el enfoque BYOD no solo es viable metodológicamente, sino que puede generar datos de alta calidad analítica sin comprometer la integridad de las mediciones, lo cual representa un avance significativo en la democratización de la investigación digital en salud. En tercer lugar, se consolidó una arquitectura computacional flexible que puede integrarse en sistemas de monitoreo remoto, lo cual constituye una contribución tecnológica con potencial de transferencia hacia el diseño de intervenciones personalizadas.

A pesar de estos avances, deben reconocerse ciertas limitaciones metodológicas y muestrales. El diseño observacional restringe la inferencia causal, de modo que futuras investigaciones longitudinales permitirán examinar la direccionalidad entre la regulación autonómica y la reducción del sedentarismo. Asimismo, aunque el protocolo BYOD favorece la representatividad ecológica, introduce heterogeneidad en la precisión de los sensores y en la adherencia de los participantes al registro continuo. Estas variaciones se mitigaron mediante limpieza automatizada de datos y validación cruzada, pero aún podrían influir marginalmente en la exactitud de las métricas. Finalmente, el tamaño muestral, aunque suficiente para el modelado difuso, podría ampliarse para explorar diferencias por sexo, edad o nivel de actividad habitual.

Las implicaciones teóricas de este estudio residen en la validación empírica de la lógica difusa como herramienta para representar la complejidad del comportamiento sedentario, superando la dicotomía tradicional entre ``activo'' e ``inactivo''. Desde el punto de vista práctico, los resultados sugieren que los modelos explicables pueden ser incorporados en aplicaciones de salud digital para generar retroalimentación personalizada basada en la respuesta fisiológica individual. Ello permitiría que las estrategias de promoción de la salud evolucionen de recomendaciones genéricas hacia intervenciones dinámicas, sensibles al contexto y orientadas a la autorregulación. En términos metodológicos, el presente trabajo ofrece un protocolo reproducible para el procesamiento, análisis y modelado de datos provenientes de wearables bajo el paradigma BYOD, alineado con los principios de ciencia abierta y acceso universal al conocimiento.

En síntesis, los hallazgos aquí presentados confirman que el monitoreo continuo del sedentarismo mediante dispositivos portátiles, interpretado a través de sistemas de inferencia difusa, representa una vía factible y ética para transformar los datos personales en conocimiento accionable para la salud pública. Este modelo no solo contribuye a la comprensión científica del sedentarismo como fenómeno biopsicosocial, sino que también proporciona una base empírica para diseñar políticas públicas orientadas a reducir la inactividad y fomentar entornos digitales saludables. Así, la investigación trasciende el ámbito académico para posicionarse como una herramienta estratégica en la promoción de estilos de vida activos y sostenibles en la era de la salud digital.

% ============================================================================
% NOTAS FINALES:
% ============================================================================
%
% RECUERDA:
% 1. INTERPRETA, no solo describas
% 2. COMPARA con estudios de tus antecedentes
% 3. EXPLICA similitudes y diferencias
% 4. Sé HONESTO con las limitaciones
% 5. Proporciona EXPLICACIONES TEÓRICAS
% 6. Conecta con implicaciones PRÁCTICAS
% 7. Verifica cumplimiento de objetivos e hipótesis
%
% ============================================================================

